% ============================================================
%  CV CONTENT. Copy this file to config.tex and edit it.
%
%    cp config.example.tex config.tex
%
%  config.tex is gitignored, so your personal details never leave
%  your machine by accident. See the Privacy section of README.md
%  if you want to version your own CV in a private repository.
%
%  main.tex and cover_letter.tex are layout only. Everything a
%  human reads is in this file, which is the only file you edit
%  per job application.
% ============================================================

% ---- Personal ----
\newcommand{\cvName}{Your Name}
\newcommand{\cvTitle}{Your Job Title}
\newcommand{\cvEmail}{you@example.com}
\newcommand{\cvPhone}{+00 000 0000}
\newcommand{\cvLocation}{City, Country}
\newcommand{\cvGithub}{your-github-username}
\newcommand{\cvLinkedin}{your-linkedin-slug}
\newcommand{\cvPortfolio}{example.com}

% ---- Summary (tailor per job) ----
% Three or four lines. Mirror the language of the job post, but only
% claim what the rest of the CV actually backs up.
\newcommand{\cvSummary}{%
One paragraph on who you are professionally: years of experience, the stack you
work in, the domain you work in, and anything you are currently studying. Keep it
to three or four lines so the first screen of the page is not all prose.}

% ---- Experience (tailor per job) ----
% \role{Title}{Dates}{Organisation}{Location}
% \sub{Heading} starts a sub-section inside a role, for a project or a team.
\newcommand{\cvExperience}{%
  \role{Your Job Title}{Mon 20XX to Present}{Company Name}{City, Country}
  {\small Previous title at the same company, Mon 20XX to Mon 20XX. One short
   clause of context if the company is not well known.}\par
  \sub{Team or Product Name (Key, Technologies, Here)}
  \begin{itemize}
    \item What you build and who uses it. Lead with the verb, name the technology,
          and end with the outcome rather than the task.
    \item A second bullet covering a different kind of work: an API, a data model,
          a migration, a piece of production support.
    \item Keep each bullet to two lines. Anything longer stops being scanned.
  \end{itemize}
  \sub{Second Team or Side of the Role (Other, Technologies)}
  \begin{itemize}
    \item Something you automated or improved, with the before and after made clear.
    \item Something you owned end to end, including the part other people avoided.
  \end{itemize}

  \role{Earlier Job Title}{Mon 20XX to Mon 20XX}{Earlier Company Name}{City, Country}
  \begin{itemize}
    \item What you shipped there, in one or two bullets. Older roles get less space.
    \item Any award or recognition, with the awarding body named.
  \end{itemize}
}

% ---- Projects (tailor per job; delete the section if you have none) ----
\newcommand{\cvProjects}{%
  \role{Project Name: What It Actually Does}{20XX to Present}{Context, for example MSc Research}{Location}
  \begin{itemize}
    \item The problem and your approach, naming the main technologies.
    \item The interesting technical decision, in plain language. No jargon a
          non-specialist reader cannot follow.
    \item How you evaluated it, with plain numbers rather than statistics jargon.
  \end{itemize}
}

% ---- Education ----
\newcommand{\cvEducation}{%
  \textbf{Your Degree, Specialisation}, \textit{\href{https://example.edu}{University Name}} \hfill {\color{muted}Mon 20XX to Present}\\
  \textbf{Your Earlier Degree}, \textit{\href{https://example.edu}{University Name}} \hfill {\color{muted}20XX to 20XX}\\
  Department Name | Classification, GPA: X.XX / X.XX
}

% ---- Skills (tailor per job) ----
% Order each line so the most relevant item for this job comes first.
\newcommand{\cvSkillLang}{Language, Language, Language}
\newcommand{\cvSkillFw}{Framework, Framework, Framework}
\newcommand{\cvSkillDb}{Database, Database}
\newcommand{\cvSkillOps}{Tool, Tool, Tool, Tool}
\newcommand{\cvSkillOther}{Concept, Concept, Concept}

% ---- References ----
% Ask permission before listing anyone, and do not put their phone number in a
% file you might ever make public. Set both to empty to drop the section.
\newcommand{\cvRefereeA}{%
  \textbf{Referee Name}\\
  Their Job Title\\
  Their Organisation\\
  their.email@example.com
}
\newcommand{\cvRefereeB}{%
  \textbf{Second Referee Name}\\
  Their Job Title\\
  Their Organisation\\
  their.email@example.com
}

% ---- Cover letter (tailor per job) ----
\newcommand{\cvCoverJobTitle}{Job Title}
\newcommand{\cvCoverCompany}{Company Name}
\newcommand{\cvCoverCity}{City, Country}
\newcommand{\cvCoverRef}{Reference: N/A}
\newcommand{\cvCoverBody}{%
I am writing to apply for the \cvCoverJobTitle{} role at \cvCoverCompany. Open with
one sentence on who you are: your current role, the kind of work you do, and anything
you are studying.

Tailor this paragraph per job. Pick two or three items from your CV's Experience and
Projects sections that most closely match the job post, and explain the connection in
your own words rather than repeating the bullets.

I would welcome the opportunity to discuss how my background fits this role. Thank you
for considering my application.}
